%macbook 
\loai{Số vân sáng tại M khi biết rõ toạ độ}
\begin{ex}%[3P5K6-3][Loai1] 
Trong thí nghiệm giao thoa ánh sáng Y-âng, khoảng các giữa hai khe là 1 mm, khoảng cách từ hai khe đến màn là 1 m. Nguồn sáng S phát ánh sáng trắng có bước sóng từ $0,4\ \mu m$ đến $0,75\ \mu m$. Tại điểm M cách vân sáng trung tâm 4 mm có mấy bức xạ cho vân sáng?
\choice
{6}
{\True 5}
{7}
{4}
\loigiai{
\begin{itemize}
\item Bước sóng của bức xạ cho vân sáng tại vị trí x: $x=k.\dfrac{\lambda D}{a}\Rightarrow \lambda =\dfrac{ax}{k.D}=\dfrac{1.4}{k.1}=\dfrac{4}{k}\ \mu m$.
\item Cho $\lambda $ vào điều kiện bước sóng của ánh sáng trắng:
$\lambda _d\le \lambda \le \lambda _t\Rightarrow 0,4\le \dfrac{4}{k}\le 0,75\Rightarrow 5,3\le k\le 10$.
\item Mà k nhận các giá trị nguyên nên: $\Rightarrow k=\left\{{6,\ 7,\ 8,\ 9,\ 10}\right\}\Rightarrow$ Có 5 bức xạ có vân sáng tại M.
\end{itemize}
}%<MyLT1>
\end{ex} \haidong{20}
\begin{ex}%[3P5K6-3][Loai1]
Trong thí nghiệm giao thoa Y-âng đối với ánh sáng trắng khoảng cách từ 2 nguồn đến màn là 2 m, khoảng cách giữa 2 nguồn là 2 mm. Số bức xạ cho vân sáng tại M cách vân trung tâm 4 mm là
\choice
{4}
{7}
{6}
{\True 5}
\loigiai{
\begin{itemize}
\item Tại điểm M có vân sáng khi đó ${x_M}=k\dfrac{\lambda D}{a}\Rightarrow \lambda =\dfrac{a.x_M}{kD}=\dfrac{2. 4}{k. 2}=\dfrac{4}{k} $.
\item Mà $0,4\le \lambda \le 0,7\Rightarrow 5,7\le k\le 10\Rightarrow k\in \left\{6;7;8;9;10 \right\} $.
\item Tại M có 5 bức xạ cho vân sáng.
\end{itemize}
}%<MyLT1>
\end{ex} \haidong{20}
\loai{Số vân sáng tại M khi phải tìm toạ độ}
\begin{ex}%[3P5G6-3][Loai2]
Trong thí nghiệm Y-âng về giao thoa ánh sáng, hai khe được chiếu bằng ánh sáng trắng có bước sóng từ 380 nm đến 760 nm. Khoảng cách giữa hai khe là 0,8\ mm, khoảng cách từ mặt phẳng chứa hai khe đến màn quan sát là 2\ m. Tại vị trí vân sáng bậc 4 của bức xạ đơn sắc có bước sóng 760 nm còn có bao nhiêu vân sáng của các bức xạ đơn sắc khác?
\choice
{5}
{2}
{3}
{\True 4}
\loigiai{
\begin{itemize}
\item Tại M là vân sáng bậc 4 của bức xạ đơn sắc \bth{760\ nm=0,76\microm} thì
\begin{equation}\rm
x_M=4\dfrac{0,76.D}{a}\tag{1}.
\end{equation}
\item Tại M còn có các vân sáng của các bức xạ khác nên
\begin{equation}\rm
x_M=k\dfrac{\lambda.D}{a}\tag{2}.
\end{equation}
\item Từ $(1)$ và $(2)$ ta có: \bth{4\dfrac{0,76.D}{a}=k\dfrac{\lambda.D}{a}\Rightarrow\lambda=\dfrac{4.0,76}{k}=\dfrac{3,04}{k}}.
\item Mặt khác: \bth{380\ nm\leq\lambda\leq760\ nm\Leftrightarrow 0,38\microm\leq\lambda\leq 0,76\microm} nên
\begin{align*}
0,38\leq \dfrac{3,04}{k}\leq 0,76\Rightarrow \dfrac{3,04}{0,76}\leq k\leq\dfrac{3,04}{0,38}\Rightarrow 4\leq k\leq 8\Rightarrow k=4;\ldots;8.
\end{align*}
%\item Có 5 giá trị nguyên của k điều này có nghĩa là có 5 bức xạ đơn sắc cho vân sáng tại M. 
%\item Do vân sáng bậc 4 là 1 trong 5 bức xạ trên nên chỉ có 4 bức xạ khác cho vân sáng tại M. 
\end{itemize}
}
%<MyLT>
\end{ex} \haidong{20}
\loai{Tìm chính xác bức xạ cho vân sáng tại M}
\begin{ex}%[3P5G6-3][Loai3]
Hai khe Y-âng cách nhau 1 mm được chiếu bằng ánh sáng trắng có bước sóng từ $0,4\ \mu m $ đến $0,7\ \mu m $, khoảng cách từ hai khe đến màn là 1 m. Tại điểm A trên màn cách vân trung tâm 2 mm có các bức xạ cho vân sáng có bước sóng
\choice
{\True $0,40\ \mu m;\ 0,50\ \mu m $ và $0,66\ \mu m $}
{ $0,44\ \mu m;\ 0,50\ \mu m $ và $0,66\ \mu m $}
{ $0,40\ \mu m;\ 0,44\ \mu m $ và $0,50\ \mu m $}
{ $0,40\ \mu m;\ 0,44\ \mu m $ và $0,66\ \mu m $}
\loigiai{
\begin{itemize}
\item Tại điểm A có vân sáng khi đó ${x_A}=k\dfrac{\lambda D}{a}\Rightarrow \lambda =\dfrac{a.x_A}{kD}=\dfrac{1. 2}{k. 1}=\dfrac{2}{k} $.
\item Mà $0,4\le \lambda \le 0,7\Rightarrow 2,86\le k\le 5\Rightarrow k\in \left\{3;4;5 \right\}\Rightarrow \lambda = \left\{0,66;0,5;0,4 \right\} $.
\item 
\item 
\end{itemize}
}
\end{ex} \haidong{20}

\begin{ex}%[3P5G6-3][Loai3]
Thực hiện giao thoa ánh sáng qua khe Y-âng, biết khoảng cách hai khe hẹp 1 mm, khoảng cách hai
khe tới màn là 2 m. Nguồn phát ánh sáng gồm các bức xạ đơn sắc có bước sóng từ $0,4\ \mu m$ đến $0,76\ \mu m$. Bước sóng lớn nhất trong số các bức xạ cho vân tối tại điểm M,
cách vân trung tâm 12 mm là
\choice
{ $0,685\ \mu m $}
{ $735\ \mu m $}
{ $0,635\ \mu m $}
{\True $0,706\ \mu m $}
\loigiai{
\begin{itemize}
\item Ta có: ${x_M}=k\dfrac{\lambda D}{a}\Rightarrow \lambda =\dfrac{a.x_M}{kD}=\dfrac{6}{k} $ (với k bán nguyên).
\item Mà $0,4\le \lambda \le 0,76\Rightarrow 7,9\le k\le 15\Rightarrow k\in \left\{8,5;9,5;10,5;11,5;12,5;13,5;14,5 \right\} \Rightarrow {{k}_{\min}}=8,5 \Rightarrow {{\lambda}_{\max}}\approx 0,706\ \mu m$.
\end{itemize}
}
\end{ex} \haidong{20}
\begin{ex}%[3P5G6-3][Loai3]
Thực hiện giao thoa ánh sáng qua khe Y-âng, biết khoảng cách hai khe hẹp 0,1 mm, khoảng cách hai
khe tới màn là 30 cm. Nguồn phát ánh sáng gồm các bức xạ đơn sắc có bước sóng từ $0,38\ \mu m $ đến $0,76\ \mu m $. Trên màn, tại vị trí vân tối thứ ba (tính từ vân trung tâm) của bức xạ $0,5\ \mu m $ có vân sáng của các bức xạ
\choice
{ $0,58\ \mu m $ và $0,44\ \mu m $}
{\True $0,42\ \mu m $ và $0,625\ \mu m $}
{ $0,625\ \mu m $}
{ $0,58\ \mu m $}
\loigiai{
\begin{itemize}
\item Theo bài ra $x=2,5\dfrac{0,5. 0,3}{0,1}=3,75\ mm $.
\item Tại điểm đó có vân sáng khi: $x=k\dfrac{\lambda D}{a}\Rightarrow \lambda =\dfrac{a.x}{kD}=\dfrac{0,1. 3,75}{k. 0,3}=\dfrac{1,25}{k} $.
\item Mà $0,38\le \lambda \le 0,76\Rightarrow 1,64\le k\le 3,29\Rightarrow k\in \left\{2;3 \right\}\Rightarrow \lambda \in \left\{0,625;0,42 \right\} $.
\item 
\item 
\end{itemize}
}
\end{ex} \haidong{20}

\loai{Số vân tối tại M}
\begin{ex}%[3P5K6-3][Loai4]
Thực hiện giao thoa ánh sáng qua khe Y-âng, biết khoảng cách hai khe hẹp 2 mm, khoảng cách hai
khe tới màn là 2 m. Nguồn phát ánh sáng gồm các bức xạ đơn sắc có bước sóng từ $0,4\ \mu m $ đến $0,76\ \mu m $. Trên màn, tại điểm cách vân trung tâm 3,3 mm có bao nhiêu bức xạ cho vân tối?
\choice
{6}
{3}
{\True 4}
{5}
\loigiai{
\begin{itemize}
\item Ta có: ${x_M}=k\dfrac{\lambda D}{a}\Rightarrow \lambda =\dfrac{a.x_M}{kD}=\dfrac{3,3}{k} $ (với k bán nguyên).
\item Mà $0,4\le \lambda \le 0,76\Rightarrow 4,34\le k\le 8,25\Rightarrow k\in \left\{4,5;5,5;6,5;7,5 \right\} $.
\item 
\item 
\end{itemize}
}%<MyLT1>
\end{ex} \haidong{20}
%\begin{ex}%[3P5G6-3][Loai4]
%Thực hiện giao thoa ánh sáng với thiết bị của Y-âng, khoảng cách giữa hai khe $a=2\ mm$, từ hai khe đến màn $D=5\ m $. Người ta chiếu sáng hai khe bằng ánh sáng trắng ($0,4\ \mu m\leq \lambda \leq 0,75\ \mu m$). Quan sát điểm A trên màn ảnh, cách vân sáng trung tâm 3,3\ mm. Tại A bức xạ cho vân tối có bước sóng ngắn nhất bằng
%\choice
%{$0,440\ \mu m$}
%{\True $0,528\ \mu m$}
%{$0,400\ \mu m$}
%{$0,490\ \mu m$}
%\loigiai{ 
%\begin{itemize}
%\item Ta có: $x_M=\left({m+0,5}\right)\dfrac{\lambda D}{a} \Rightarrow \lambda =\dfrac{ax_M}{\left({m+0,5}\right)D}=\dfrac{3,3}{m+0,5}\ \mu m\xrightarrow{{0,4\leq \lambda =\dfrac{3,3}{m+0,5}\leq 0,75}} 3,9\leq m\leq 7,75\Rightarrow m=4;5;6;7\\
%\Rightarrow \lambda_{\min}=\dfrac{3,3}{7+0,5}=0,44\ \mu m$.
%\item 
%\item 
%\end{itemize}
%}
%\end{ex} \haidong{20}
\begin{ex}%[3P5G6-3][Loai4]
Thực hiện giao thoa ánh sáng với thiết bị của Y-âng, khoảng cách giữa hai khe 2 mm, từ hai khe đến màn 2 m. Người ta chiếu sáng hai khi bằng ánh sáng trắng có bước sóng từ $0,38\ \mu m $ đến $0,76\ \mu m $. Quan sát điểm M trên màn ảnh, cách vân sáng trung tâm 3,3 mm. Tại M bức xạ cho vân tối có bước sóng ngắn nhất bằng
\choice
{490 nm}
{508 nm}
{\True 388 nm}
{440 nm}
\loigiai{
\begin{itemize}
\item Ta có: ${x_M}=k\dfrac{\lambda D}{a}\Rightarrow \lambda =\dfrac{a.x_M}{kD}=\dfrac{3,3}{k} $ (với k bán nguyên).
\item Mà $0,38\le \lambda \le 0,76\Rightarrow 4,34\le k\le 8,68\Rightarrow k\in \left\{4,5;5,5;6,5;7,5;8,5 \right\} $ $\Rightarrow {{k}_{\max}}=8,5 \Rightarrow {{\lambda}_{\min}}\approx 0,388\ \mu m $.
\item 
\item 
\end{itemize}
}
\end{ex} \haidong{20}
\loai{Bức xạ không cho vân giao thoa}
\begin{ex}%[3P5G6-3][Loai5]
Trong thí nghiệm giao thoa Y-âng, khoảng cách giữa hai khe là 0,9 mm, khoảng cách từ mặt phẳng hai khe đến màn là 1 m. Khe S được chiếu bằng ánh sáng trắng có bước sóng từ $0,38\ \mu m $ đến $0,76\ \mu m $. Bức xạ đơn sắc nào sau đây \textbf{không} cho vân sáng tại điểm cách vân trung tâm 3 mm?
\choice
{450 nm}
{\True 650 nm}
{540 nm}
{675 nm}
\loigiai{
\begin{itemize}
\item Tại điểm A có vân sáng khi đó ${x_A}=k\dfrac{\lambda D}{a}\Rightarrow \lambda =\dfrac{a.x_A}{kD}=\dfrac{0,9. 3}{k. 1}=\dfrac{27}{10k} $.
\item Mà $0,38\le \lambda \le 0,76\Rightarrow 3,55\le k\le 7,1\Rightarrow k\in \left\{4;5;6;7 \right\}\Rightarrow \lambda \in \left\{0,675;0,54;0,45;\dfrac{27}{70} \right\} $.
\item 
\item 
\end{itemize}
}
\end{ex} \haidong{20}

